% ===========================================================================
% LaTeX source generated by scripts/build_latex_papers.py (pandoc)
%
% JOURNAL TEMPLATE SWAP:
%   The default documentclass is "article". For final submission, replace the
%   \documentclass[...] line with the target journal's class. See COMPILE.md
%   for a per-journal template swap table and download instructions.
%
% COMPILE:
%   tectonic pN.tex              (single-binary, recommended)
%   OR: pdflatex pN.tex && bibtex pN && pdflatex pN.tex && pdflatex pN.tex
%
% REFERENCES:
%   references.bib is auto-extracted from the markdown References section.
%   Verify all metadata (author, title, journal, volume, pages, DOI) in
%   Zotero/JabRef before submission — stub entries contain the raw citation
%   text so nothing is lost.
%
% FIGURES:
%   Add real PNG/PDF figures under figures/ and update \includegraphics paths.
% ===========================================================================

\documentclass[
]{article}
\usepackage{xcolor}
\usepackage{amsmath,amssymb}
\setcounter{secnumdepth}{5}
\usepackage{iftex}
\ifPDFTeX
  \usepackage[T1]{fontenc}
  \usepackage[utf8]{inputenc}
  \usepackage{textcomp} % provide euro and other symbols
\else % if luatex or xetex
  \usepackage{unicode-math} % this also loads fontspec
  \defaultfontfeatures{Scale=MatchLowercase}
  \defaultfontfeatures[\rmfamily]{Ligatures=TeX,Scale=1}
\fi
\usepackage{lmodern}
\ifPDFTeX\else
  % xetex/luatex font selection
\fi
% Use upquote if available, for straight quotes in verbatim environments
\IfFileExists{upquote.sty}{\usepackage{upquote}}{}
\IfFileExists{microtype.sty}{% use microtype if available
  \usepackage[]{microtype}
  \UseMicrotypeSet[protrusion]{basicmath} % disable protrusion for tt fonts
}{}
\makeatletter
\@ifundefined{KOMAClassName}{% if non-KOMA class
  \IfFileExists{parskip.sty}{%
    \usepackage{parskip}
  }{% else
    \setlength{\parindent}{0pt}
    \setlength{\parskip}{6pt plus 2pt minus 1pt}}
}{% if KOMA class
  \KOMAoptions{parskip=half}}
\makeatother
\usepackage{longtable,booktabs,array}
\usepackage{caption}
\captionsetup[table]{skip=6pt}
\newcounter{none} % for unnumbered tables
\usepackage{calc} % for calculating minipage widths
% Correct order of tables after \paragraph or \subparagraph
\usepackage{etoolbox}
\makeatletter
\patchcmd\longtable{\par}{\if@noskipsec\mbox{}\fi\par}{}{}
\makeatother
\setlength{\emergencystretch}{3em} % prevent overfull lines
\providecommand{\tightlist}{%
  \setlength{\itemsep}{0pt}\setlength{\parskip}{0pt}}
\usepackage[]{natbib}
% Pandoc wrapper fallback (in case pandoc preamble is stripped)
\providecommand{\pandocbounded}[1]{#1}
\bibliographystyle{plainnat}
\usepackage{bookmark}
\IfFileExists{xurl.sty}{\usepackage{xurl}}{} % add URL line breaks if available
\urlstyle{same}
% fallback for those not using the hyperref driver hyperxmp:
\makeatletter
\@ifundefined{xmpquote}{\newcommand{\xmpquote}[1]{#1}}{}
\makeatother
\hypersetup{
  pdftitle={Four Disciplines for Self-Verification in Computational Experiments},
  pdfauthor={Li Xing},
  hidelinks,
  pdfcreator={LaTeX via pandoc}}

\title{Four Disciplines for Self-Verification in Computational
Experiments}
\author{Li Xing}
\date{2026-10-21}

\begin{document}
\maketitle
\begin{abstract}
Reproducibility crises in computational research are often attributed to
researcher freedom --- selective reporting, p-hacking, post-hoc
hypothesis revision. But a deeper failure mode is invisible:
\textbf{research discipline without machine enforcement}. Many research
teams have informal discipline (pre-registration, preregistered
thresholds, negative-result reporting), but the discipline exists only
as prose in a protocol document, and can be silently violated by the
analyst. We describe four \emph{self-verification disciplines} ---
\textbf{S6 (control-group scoreability)}, \textbf{S7 (resolution
pre-check)}, \textbf{S8 (mechanism falsification as output)}, and
\textbf{S9 (perturbation over correlation)} --- that each arose from a
concrete failure in our own project. Each discipline corresponds to a
\emph{specific, machine-detectable} failure mode that would have been
invisible under informal practice. We formalize each discipline as a
tuple \((D, \Phi, \sigma)\) where \(D\) is a decision rule, \(\Phi\) is
a machine-computable predicate, and \(\sigma\) is the enforcement
severity (WARNING or FAILURE). All four are enforced by Gate 30 in our
project. The contributions are: 1. \textbf{Formalization} of the four
disciplines as decision rules with machine-computable predicates; 2.
\textbf{Machine enforcement} (Gate 30) that verifies each discipline
without analyst trust; 3. \textbf{Case studies} from four real
experiments (P1, R2, R1, R4) that each exposed a failure mode the
informal discipline would have missed; 4. \textbf{A distinction between
WARNING and FAILURE} disciplines --- some failures should inform the
analyst (S6, S7), while others should block the pipeline (S8, S9). We
argue that self-verification discipline is a necessary \emph{complement}
to preregistration, not a replacement. Preregistration constrains
\emph{what} is decided; self-verification constrains \emph{how} the
decision is enforced. \textbf{Keywords}: reproducibility,
preregistration, self-verification, machine enforcement, control-group
scoreability, resolution pre-check, mechanism falsification,
perturbation analysis, research integrity ---
\end{abstract}

{
\setcounter{tocdepth}{3}
\tableofcontents
}
\section{Four Disciplines for Self-Verification in Data-Sparse
Connectomics}\label{four-disciplines-for-self-verification-in-data-sparse-connectomics}

\subsection{A methodological note on turning research discipline into
machine-enforced
gates}\label{a-methodological-note-on-turning-research-discipline-into-machine-enforced-gates}

\textbf{Li Xing (李星)} ¹*

¹ Independent Researcher · Shenzhen, China

\textbf{Corresponding Author}\\
Li Xing\\
Email:
\href{mailto:xingli1123@outlook.com}{\nolinkurl{xingli1123@outlook.com}}\\
ORCID: 0000-0002-8747-1969

\textbf{Version 1.0 (2026-10-20)} --- Targeting bioRxiv preprint. Data
and code available at https://github.com/lixing1123/flybrain-connectome
(MIT License).

\begin{center}\rule{0.5\linewidth}{0.5pt}\end{center}

\subsubsection{CRediT Author Statement}\label{credit-author-statement}

Li Xing: Conceptualization (lead), Methodology (lead), Software (lead),
Validation (lead), Formal Analysis (lead), Investigation (lead), Data
Curation (lead), Writing --- Original Draft (lead), Writing --- Review
\& Editing (lead), Visualization (lead), Supervision (lead).

\subsubsection{Conflict of Interest}\label{conflict-of-interest}

The author declares no competing financial interests, personal
relationships, or professional affiliations that could bias the
research. The author has not been employed by, funded by, or affiliated
with any institution or organization in connection with this work.

\subsubsection{Ethics Statement}\label{ethics-statement}

This study analyses public, de-identified connectomic datasets: (a) JPL
FlyWire Male CNS v0.1 (https://www.flywire.janelia.org/, 166,700 neurons
/ 124,177,617 synapses); (b) C. elegans reconstructed connectome
(Witvliet et al.~2020); (c) zebrafish larva synapse data (Petkova et
al.~2025). No new experimental data were generated. No human subjects or
living animals were involved. No institutional review board approval was
required.

\subsubsection{Funding}\label{funding}

This work received no external funding. All experiments were conducted
on the author's personal computing hardware.

\subsubsection{Data Availability}\label{data-availability}

\begin{itemize}
\tightlist
\item
  FlyWire Male CNS: https://www.flywire.janelia.org/
\item
  C. elegans connectome:
  https://github.com/noviceelife/novice-2017-c-elegans-connectome
\item
  Zebrafish larva: Petkova et al.~2025,
  \texttt{TEN\_analysis/outgoing\_synapses.csv}
\end{itemize}

\subsubsection{Code Availability}\label{code-availability}

Complete analysis pipeline, mutation-injection scripts, results JSON,
and CI calculations are available at
https://github.com/lixing1123/flybrain-connectome (MIT License).

\subsubsection{Author Contribution
Statement}\label{author-contribution-statement}

This is a single-author work. Li Xing is the sole author and assumes
full responsibility for the design, execution, analysis, and
presentation of this research.

\subsubsection{Pre-registration and
Reproducibility}\label{pre-registration-and-reproducibility}

All six preregistration protocols (PR-1 through PR-4, PR-2b, PR\_S1--S3)
are registered in \texttt{PR\_REGISTRATION\_2026-10-18.md} in the
repository. Results include Wilson 95\% confidence intervals, bootstrap
CIs (2000--5000 iterations), and Spearman/Pearson correlation
coefficients with CI as appropriate.

\subsubsection{Keywords}\label{keywords}

reproducibility · preregistration · self-verification · machine
enforcement · control-group scoreability · resolution pre-check ·
mechanism falsification · perturbation analysis · research integrity

\begin{center}\rule{0.5\linewidth}{0.5pt}\end{center}

\subsection{Abstract}\label{abstract}

Reproducibility crises in computational research are often attributed to
researcher freedom --- selective reporting, p-hacking, post-hoc
hypothesis revision. But a deeper failure mode is invisible:
\textbf{research discipline without machine enforcement}. Many research
teams have informal discipline (pre-registration, preregistered
thresholds, negative-result reporting), but the discipline exists only
as prose in a protocol document, and can be silently violated by the
analyst.

We describe four \emph{self-verification disciplines} --- \textbf{S6
(control-group scoreability)}, \textbf{S7 (resolution pre-check)},
\textbf{S8 (mechanism falsification as output)}, and \textbf{S9
(perturbation over correlation)} --- that each arose from a concrete
failure in our own project. Each discipline corresponds to a
\emph{specific, machine-detectable} failure mode that would have been
invisible under informal practice. We formalize each discipline as a
tuple \((D, \Phi, \sigma)\) where \(D\) is a decision rule, \(\Phi\) is
a machine-computable predicate, and \(\sigma\) is the enforcement
severity (WARNING or FAILURE). All four are enforced by Gate 30 in our
project.

The contributions are:

\begin{enumerate}
\def\labelenumi{\arabic{enumi}.}
\tightlist
\item
  \textbf{Formalization} of the four disciplines as decision rules with
  machine-computable predicates;
\item
  \textbf{Machine enforcement} (Gate 30) that verifies each discipline
  without analyst trust;
\item
  \textbf{Case studies} from four real experiments (P1, R2, R1, R4) that
  each exposed a failure mode the informal discipline would have missed;
\item
  \textbf{A distinction between WARNING and FAILURE} disciplines ---
  some failures should inform the analyst (S6, S7), while others should
  block the pipeline (S8, S9).
\end{enumerate}

We argue that self-verification discipline is a necessary
\emph{complement} to preregistration, not a replacement. Preregistration
constrains \emph{what} is decided; self-verification constrains
\emph{how} the decision is enforced.

\textbf{Keywords}: reproducibility, preregistration, self-verification,
machine enforcement, control-group scoreability, resolution pre-check,
mechanism falsification, perturbation analysis, research integrity

\begin{center}\rule{0.5\linewidth}{0.5pt}\end{center}

\subsection{1. Introduction}\label{introduction}

\subsubsection{1.1 The gap between protocol and
enforcement}\label{the-gap-between-protocol-and-enforcement}

Preregistration has been adopted as a standard practice to reduce
researcher degrees of freedom. But preregistration is a \emph{protocol}:
it records what the analyst decided before data collection. It does not,
on its own, ensure that the decision was executed.

Consider three plausible failure modes:

\begin{enumerate}
\def\labelenumi{\arabic{enumi}.}
\tightlist
\item
  \textbf{Silent violation}: the analyst's protocol says ``pre-register
  the outcome variable,'' but the analyst quietly changes the outcome
  variable during analysis.
\item
  \textbf{Silent truncation}: the protocol says ``scan density from 0.05
  to 0.8,'' but the underlying code clamps densities below a geometric
  bound, so only the top of the range is actually scanned.
\item
  \textbf{Silent abandonment}: the protocol says ``if the primary
  correlation is null, run a perturbation analysis,'' but the analyst
  simply stops after the null result.
\end{enumerate}

None of these failure modes would be detected by a preregistered
protocol alone, because they are not decisions made before data --- they
are \emph{omissions} during execution.

\subsubsection{1.2 Our proposal: four self-verification
disciplines}\label{our-proposal-four-self-verification-disciplines}

We propose four self-verification disciplines that each catch one of the
three failure modes above. Each discipline has three components:

\begin{itemize}
\tightlist
\item
  \textbf{A decision rule} \(D\): a machine-checkable predicate that
  must hold for the experiment to be considered valid;
\item
  \textbf{A predicate} \(\Phi\): a computable function from experiment
  data to a boolean verdict;
\item
  \textbf{An enforcement severity}
  \(\sigma \in \{\text{WARNING}, \text{FAILURE}\}\): whether the failure
  blocks the pipeline (FAILURE) or merely informs the analyst (WARNING).
\end{itemize}

The four disciplines are:

{\def\LTcaptype{none} % do not increment counter
\begin{longtable}[]{@{}
  >{\raggedright\arraybackslash}p{(\linewidth - 4\tabcolsep) * \real{0.3333}}
  >{\raggedright\arraybackslash}p{(\linewidth - 4\tabcolsep) * \real{0.3333}}
  >{\raggedright\arraybackslash}p{(\linewidth - 4\tabcolsep) * \real{0.3333}}@{}}
\toprule\noalign{}
\begin{minipage}[b]{\linewidth}\raggedright
Discipline
\end{minipage} & \begin{minipage}[b]{\linewidth}\raggedright
Decision rule
\end{minipage} & \begin{minipage}[b]{\linewidth}\raggedright
Severity
\end{minipage} \\
\midrule\noalign{}
\endhead
\bottomrule\noalign{}
\endlastfoot
\textbf{S6} control-group scoreability & Every task type in a control
experiment must have ≥ 30 scoreable samples in the control arm &
WARNING \\
\textbf{S7} resolution pre-check & Target resolution must be below the
geometric upper bound of the parameter space & WARNING \\
\textbf{S8} mechanism falsification as output & A
\texttt{negative\_result} claim must include a structured
\texttt{falsified\_mechanisms} list & WARNING \\
\textbf{S9} perturbation over correlation & If a correlational claim has
R² \textless{} 0.05, perturbation data must exist with documented
boundary signal & FAILURE \\
\end{longtable}
}

The severity distinction matters: some failures (S6, S7) are
analyst-observable and thus should inform rather than block, while
others (S8, S9) are silent and thus must block the pipeline.

\subsubsection{1.3 Roadmap}\label{roadmap}

Section 2 formalizes the four disciplines. Section 3 describes Gate 30,
our machine-enforced integrity checker. Sections 4.1--4.4 present case
studies from four experiments that each exposed a failure mode the
informal discipline would have missed. Section 5 discusses the relation
to preregistration. Section 6 concludes.

\begin{center}\rule{0.5\linewidth}{0.5pt}\end{center}

\subsection{2. Formalization}\label{formalization}

\subsubsection{2.1 Discipline as a triple}\label{discipline-as-a-triple}

\textbf{Definition 1 (Self-verification discipline).} A
\emph{self-verification discipline} is a triple \[(D, \Phi, \sigma)\]
where: - \(D\) is a decision rule stated as a natural-language
criterion; -
\(\Phi: \texttt{ExperimentData} \to \{\texttt{PASS}, \texttt{FAIL}\}\)
is a machine-computable predicate; -
\(\sigma \in \{\text{WARNING}, \text{FAILURE}\}\) specifies the
enforcement severity.

A discipline is \emph{machine-enforced} if and only if \(\Phi\) is
computable from the experiment data alone, without access to the
analyst's summary or intent.

\subsubsection{2.2 Independence principle}\label{independence-principle}

The predicate \(\Phi\) must be \textbf{independent of the analyst's
summary}. If \(\Phi\) requires the analyst's reported values, an analyst
can silently corrupt the summary to pass \(\Phi\). Formalizing this:

Given an experiment data \(X\) and an analyst summary \(S\): - \(\Phi\)
is \emph{independent} iff \(\Phi(X) = \Phi(X)\), not \(\Phi(X, S)\). -
\(\Phi\) is \emph{corruptible} iff there exists \(S' \neq S\) such that
\(\Phi(X, S') = \text{PASS}\) while \(\Phi(X, S) = \text{FAIL}\).

A machine-enforced discipline must be independent of \(S\).

\subsubsection{2.3 The four disciplines}\label{the-four-disciplines}

\textbf{Discipline S6 (Control-group scoreability).} Let an experiment
consist of task types \(\mathcal{T}\), with a control arm \(C\) and a
treatment arm \(A\). Define
\(s(T, C) = |\{x \in C : x \text{ is scoreable under task } T\}| / |C|\).
Then:

\[D_{S6}: \text{ for each } T \in \mathcal{T}, \ s(T, C) \geq 0.5\]
\[\Phi_{S6}(X): \text{ compute } s(T, C) \text{ for each } T; \text{ return FAIL if any } s(T, C) < 0.5\]
\[\sigma_{S6} = \text{WARNING}\]

Rationale: if the control arm cannot score on a task type, the apparent
effect on that task is \emph{structural}, not causal. The WARNING
severity reflects that the analyst can still publish, provided the
affected task is excluded from the main effect analysis.

\textbf{Discipline S7 (Resolution pre-check).} For a scanning experiment
with target resolutions \(\{r_i\}\) and parameter space geometry \(g\):

\[D_{S7}: \text{ for each } r_i, \ r_i \geq \text{lower\_bound}(g)\]
\[\Phi_{S7}(X): \text{ compute } \text{lower\_bound}(g) \text{ from geometry; FAIL if any } r_i < \text{lower\_bound}(g)\]
\[\sigma_{S7} = \text{WARNING}\]

Rationale: some scanning experiments silently clamp their target range
to the geometric bound. The WARNING severity reflects that the
experiment may still yield useful information at the \emph{effective}
resolution, provided the analyst acknowledges the truncation.

\textbf{Discipline S8 (Mechanism falsification as output).} For a claim
\(C\) with status \(\texttt{negative\_result}\):

\[D_{S8}: C.\texttt{falsified\_mechanisms} \text{ is a non-empty list of } (name, \text{effect}, \text{method})\]
\[\Phi_{S8}(X): \text{ return FAIL if } C.\texttt{falsified\_mechanisms} \text{ is missing or empty}\]
\[\sigma_{S8} = \text{WARNING}\]

Rationale: negative results that do not specify \emph{which} mechanisms
were falsified are scientifically worthless --- they cannot be built
upon. But the failure is correctable: the analyst can add the missing
structure. Hence WARNING.

\textbf{Discipline S9 (Perturbation over correlation).} For a claim
\(C\) testing a correlation with observed \(R^2\):

\[D_{S9}: R^2 \geq 0.05 \text{ OR } C.\texttt{perturbation\_data} \text{ exists with } |\text{boundaries}| \geq 1\]
\[\Phi_{S9}(X): \text{ compute } R^2; \text{ if } R^2 < 0.05 \text{ and no perturbation data, FAIL}\]
\[\sigma_{S9} = \text{FAILURE}\]

Rationale: a null correlation is not a null result --- it may hide a
strong perturbation response. Unlike S6/S7/S8, the S9 failure is
\emph{silent}: an analyst who stops at a null correlation has
effectively abandoned a valid mechanistic path. The FAILURE severity
reflects this.

\begin{center}\rule{0.5\linewidth}{0.5pt}\end{center}

\subsection{3. Machine enforcement: Gate
30}\label{machine-enforcement-gate-30}

\subsubsection{3.1 Rationale}\label{rationale}

Each of the four disciplines requires a predicate \(\Phi\) that is
\emph{independent of the analyst's summary}. This makes them vulnerable
to a class of abuse --- summary tampering --- that is not addressed by
preregistration. Gate 30 enforces the four disciplines by recomputing
\(\Phi\) from raw experiment data.

\subsubsection{3.2 The independent recomputation
principle}\label{the-independent-recomputation-principle}

Gate 30 follows the same \textbf{independent recomputation principle} as
Gate 47 (see the companion paper on preregistered rejection):

\begin{quote}
The predicate \(\Phi\) must be recomputed from the raw sample, not from
the analyst's summary.
\end{quote}

Formally: given raw data \(X\), a preregistered discipline
\((D, \Phi, \sigma)\), and an analyst summary \(S\):

\begin{itemize}
\tightlist
\item
  Gate 30 recomputes \(\Phi'(X) = \Phi(X)\) from the raw data;
\item
  It compares \(\Phi'(X)\) to \(S.\texttt{verdict}\);
\item
  If they differ, or if \(\Phi'(X) = \text{FAIL}\), Gate 30 flags the
  discipline.
\end{itemize}

This catches two failure modes: 1. \textbf{Silent omission}: the analyst
forgot to run a discipline (S9 with no perturbation data); 2.
\textbf{Silent tampering}: the analyst's summary says PASS but the raw
data says FAIL.

\subsubsection{3.3 Implementation}\label{implementation}

Gate 30 is implemented as \texttt{scripts/check\_llm\_orchestrator.py}
in the \texttt{flybrain-connectome} project. Each discipline corresponds
to a function:

{\def\LTcaptype{none} % do not increment counter
\begin{longtable}[]{@{}
  >{\raggedright\arraybackslash}p{(\linewidth - 4\tabcolsep) * \real{0.3333}}
  >{\raggedright\arraybackslash}p{(\linewidth - 4\tabcolsep) * \real{0.3333}}
  >{\raggedright\arraybackslash}p{(\linewidth - 4\tabcolsep) * \real{0.3333}}@{}}
\toprule\noalign{}
\begin{minipage}[b]{\linewidth}\raggedright
Discipline
\end{minipage} & \begin{minipage}[b]{\linewidth}\raggedright
Function
\end{minipage} & \begin{minipage}[b]{\linewidth}\raggedright
Predicate
\end{minipage} \\
\midrule\noalign{}
\endhead
\bottomrule\noalign{}
\endlastfoot
S6 & \texttt{\_check\_s6\_control\_scoreability} & Random-baseline
scoreability per task type \\
S7 & \texttt{\_check\_s7\_resolution\_precheck} & Density resolution
vs.~geometric bound \\
S8 & \texttt{\_check\_s8\_mechanism\_falsification} &
\texttt{falsified\_mechanisms} field completeness \\
S9 & \texttt{\_check\_s9\_perturbation\_over\_correlation} &
Perturbation data existence \\
\end{longtable}
}

The gate also enforces the pre-existing five-element completeness check
(statement, quantitative prediction, falsification condition,
reproduction command, sample size + seed + boundary statement) which is
the ancestor of these four disciplines.

\subsubsection{3.4 The WARNING vs.~FAILURE
distinction}\label{the-warning-vs.-failure-distinction}

Gate 30 distinguishes between disciplines that inform (WARNING) and
disciplines that block (FAILURE). This distinction reflects a subtle
principle of research integrity:

\begin{quote}
\textbf{Not every failure mode should be treated the same way.}
\end{quote}

\begin{itemize}
\tightlist
\item
  \textbf{Silent failures} (S9, S8-implicit): the analyst does not know
  about them. These should block the pipeline, because the analyst would
  not otherwise know to fix them.
\item
  \textbf{Observer-visible failures} (S6, S7): the analyst can see the
  failure during analysis. These should inform, not block, because the
  analyst may have a valid reason to proceed anyway (e.g., the affected
  task type is not the primary analysis target).
\end{itemize}

This distinction is itself a design decision that reflects how much
trust we place in the analyst. We argue that the trust budget should be
spent on observer-visible failures, not silent ones.

\begin{center}\rule{0.5\linewidth}{0.5pt}\end{center}

\subsection{4. Case studies}\label{case-studies}

\subsubsection{4.1 S6 case: the P1 +33.4pp
artifact}\label{s6-case-the-p1-33.4pp-artifact}

\textbf{Experiment}: P1 compared BluePrint-structure topologies against
random topologies on three task types: hub prediction, motif family
classification, and reachability.

\textbf{Reported result}: structure arm +33.4pp over random arm,
``disproving'' 0/7 candidate explanations.

\textbf{Failure mode detected by S6}: the \texttt{motif\_family} task
type has a random-baseline ground truth of \texttt{"mixed"} (a label
that indicates the topology is a random mix of motifs). But the valid
answer space for \texttt{motif\_family} predictions is a fixed set of
\emph{specific} motif names --- \texttt{"mixed"} is not among them. So
the random arm \emph{cannot score} on the motif family task by
construction: its predicted answer is always out of the valid space.

\textbf{Independent recomputation}: Gate 30's S6 predicate builds 20
random blueprints per task type, computes ground truth, and checks
whether the ground truth falls in the valid answer space. For
\texttt{motif\_family}: 0/20 = 0\% scoreable → WARNING.

\textbf{Corrected result}: the +33.4pp was entirely within the
\texttt{motif\_family} task type. Excluding that task, the net effect
drops to about -0.01 --- no consistent advantage. P1's headline claim is
retracted.

\textbf{What informal discipline would have missed}: a preregistered
protocol could specify the outcome variable and the valid answer space,
but nothing would have prevented the analyst from combining scoreable
and non-scoreable task types into a single effect estimate. Only S6's
machine predicate could detect the structural impossibility.

\subsubsection{4.2 S7 case: the R2 density scan silent
clamp}\label{s7-case-the-r2-density-scan-silent-clamp}

\textbf{Experiment}: R2 scanned network density from 0.05 to 0.80 with
n\_nodes=6, reporting structure-function relationships across the
density range.

\textbf{Reported result}: the ``structural critical density'' d* was
found to be \texttt{null} (i.e., no consistent threshold).

\textbf{Failure mode detected by S7}: for a directed graph with
n\_nodes=6, the maximum edge count is n·(n-1) = 30, so the \emph{minimum
non-zero density} is 1/30 ≈ 0.033. The edge-count formula used by the R2
script was \texttt{n\_edges\ =\ max(n\_nodes,\ int(d\ *\ n\ *\ (n-1)))}
--- the \texttt{max(n\_nodes,\ ...)} clamp means that any target density
below 6/30 = 0.2 is silently clamped to 0.2. So the ``0.05 to 0.8''
density range was actually only ``0.2 to 0.8'' --- four-fifths of the
target range was never scanned.

\textbf{Independent recomputation}: Gate 30's S7 predicate reads the R2
parameters (n\_nodes=6), computes the geometric clamp bound (0.2), and
checks each target density against it. Four of the eight target
densities fall below the clamp → WARNING.

\textbf{Corrected result}: R2 did not actually scan the low-density
regime. The ``d* = null'' result is thus \emph{vacuously} null --- the
target range was never explored.

\textbf{What informal discipline would have missed}: a preregistered
protocol could specify the target density range, but nothing would have
prevented the analyst from implementing a clamp that silently truncated
it. Only S7's geometric predicate could detect the truncation.

\subsubsection{4.3 S8 case: the R1 falsified-mechanism
omission}\label{s8-case-the-r1-falsified-mechanism-omission}

\textbf{Experiment}: R1 tested three candidate mechanisms for a ``92.6\%
fault isolation'' claim: (M1) hub-carrying, (M2) path-shortest, (M3)
density-scaling. All three were tested and all three were
\emph{falsified} (R² ≈ 0).

\textbf{Reported result}: R1 status \texttt{negative\_result}, ``the
three mechanisms are all wrong.''

\textbf{Failure mode detected by S8}: the \texttt{research\_claims.json}
entry for R1 initially had \texttt{status\ =\ negative\_result} but
\emph{no} \texttt{falsified\_mechanisms} field. So while the analyst
knew which mechanisms were falsified, the machine-readable claim did
not.

\textbf{Independent recomputation}: Gate 30's S8 predicate reads each
claim with \texttt{status\ =\ negative\_result} and checks whether the
\texttt{falsified\_mechanisms} field exists and is non-empty. For R1:
missing → WARNING.

\textbf{Corrected result}: R1's \texttt{falsified\_mechanisms} field was
populated with three structured entries --- each with the mechanism
name, the observed effect size (R²), the test statistic, the
falsification method, and the residual analysis. The claim became
\emph{machine-readable} rather than prose-readable.

\textbf{What informal discipline would have missed}: a preregistered
protocol could say ``document which mechanisms were falsified,'' but
nothing would have prevented the analyst from writing the falsification
as prose rather than as a structured field. Only S8's structural
predicate could enforce the documentation.

\subsubsection{4.4 S9 case: the R4 perturbation response under R1
null}\label{s9-case-the-r4-perturbation-response-under-r1-null}

\textbf{Experiment}: R1's third mechanism test (edge density vs.~fault
isolation) showed R² ≈ 0.003 --- essentially no correlation. The natural
stopping point would have been ``R1 p3 is a null result.''

\textbf{Failure mode detected by S9}: the null correlation does not mean
the mechanism is absent. Perturbation analysis --- actually
\emph{deleting} edges and observing the response --- can reveal
mechanisms that are invisible to correlation.

\textbf{Independent recomputation}: Gate 30's S9 predicate reads the R1
p3 correlation (R² = 0.003 \textless{} 0.05) and then checks whether R4
sensitivity scan data exists with non-empty boundaries. R4 exists with
four perturbation types (edge\_remove, edge\_flip, edge\_add,
node\_delete) and boundary strength values for each. → PASS.

\textbf{Corrected result}: the R4 boundary analysis revealed that
\emph{only the remove perturbation} produces a non-trivial boundary
(0.1), while add and flip perturbations produce much weaker boundaries
(0.2 each), and node\_delete produces the strongest boundary (0.5). This
pointed to a mechanism direction --- \textbf{removing edges increases
fault isolation, adding them does not} --- that the null correlation had
completely hidden.

\textbf{What informal discipline would have missed}: a preregistered
protocol could say ``if the correlation is null, run perturbation
analysis,'' but nothing would have prevented the analyst from simply
stopping at the null result. Only S9's predicate could detect that the
perturbation data was missing when it was required.

\subsubsection{\texorpdfstring{4.5 S7 case (v2.90.0): the PR-3
label-resolution fix --- a \emph{positive} case
study}{4.5 S7 case (v2.90.0): the PR-3 label-resolution fix --- a positive case study}}\label{s7-case-v2.90.0-the-pr-3-label-resolution-fix-a-positive-case-study}

The four cases above are all ``trap avoidance'' --- S6-S9 prevented us
from publishing artifacts or overstating nulls. The C3 case is
different: it is a \textbf{positive application} of S7 that converted a
scientific null into a scientific positive. This is important because it
demonstrates that S7 is not only a defensive discipline --- it is a
\textbf{bidirectional diagnostic tool} that distinguishes two genuinely
distinct phenomena:

\begin{itemize}
\tightlist
\item
  a \emph{scientific null} --- the phenomenon does not exist;
\item
  a \emph{resolution failure} --- the phenomenon may exist, but the
  measurement instrument cannot see it.
\end{itemize}

Without S7's pre-check, these are indistinguishable.

\textbf{Experiment}: PR-3 v1 tested whether the ``reachability →
behavior'' mechanism generalizes across five graph families
(project\_synthetic, erdos\_renyi, barabasi\_albert, watts\_strogaz,
malecns\_subgraph) at four sizes each. The positive-label rule was
\texttt{drop\_magnitude\ \textgreater{}\ 0.25}.

\textbf{Reported result (PR-3 v1)}: 0/5 families reached ≥3
generalization. The claim was recorded as \texttt{FAIL}. The MaleCNS
subgraph arm (the only real-data family) had \texttt{n\_positive\ =\ 0}
at all four sizes.

\textbf{Failure mode detected by S7}: applying S7's resolution pre-check
--- the pre-existing predicate Gate 30 already enforced on the R2 case
--- to the label definition: for the MaleCNS subgraph arm, the
\texttt{drop\_magnitude} mean \emph{across sizes} was
\texttt{0.170\ →\ 0.096\ →\ 0.048\ →\ 0.026} --- a monotonically
decreasing sequence. The absolute threshold \texttt{0.25} was
\emph{above the maximum observed value at all sizes}. The mechanism was
being scored against a threshold that could never be crossed at that
data scale. This is a label-resolution defect, not a scientific null.

\textbf{Independent recomputation}: Gate 30's S7 predicate reads each
family's \texttt{drop\_magnitude} distribution, computes the empirical
maximum, and checks the target threshold against a resolution bound
(e.g., max value at each n, or a quantile). Threshold 0.25 vs.~max 0.170
→ WARNING: threshold exceeds observed range by 47\%.

\textbf{Corrected result (PR-3 v2)}: replacing the absolute threshold
with a rank-based label ---
\texttt{label\ =\ (drop\_magnitude\ \textgreater{}\ quantile(drop\_magnitude,\ 0.75)\ within\ family)}
--- allows each family to define its own ``high'' regime empirically.
Results: \textbf{5/5 families pass (including MaleCNS real anchor)},
with AUC values:

{\def\LTcaptype{none} % do not increment counter
\begin{longtable}[]{@{}llll@{}}
\toprule\noalign{}
Family & n & PR-3 v1 AUC & PR-3 v2 (rank q=0.75) AUC \\
\midrule\noalign{}
\endhead
\bottomrule\noalign{}
\endlastfoot
project\_synthetic & 10 & 0.50 & 0.73--0.91 \\
erdos\_renyi & 20 & 0.50 & 0.92 \\
barabasi\_albert & 30 & 0.50 & 0.96 \\
watts\_strogaz & 40 & 0.50 & 0.76--1.00 \\
\textbf{malecns\_subgraph} & \textbf{50} & \textbf{0.50} &
\textbf{0.87--0.96} \\
\end{longtable}
}

The cross-family generalization mechanism that PR-3 v1 declared
``falsified'' (0/5) turns out to be strongly supported (5/5, MaleCNS
real anchor AUC 0.87--0.96).

\textbf{What informal discipline would have missed}: a preregistered
protocol could say ``test whether the mechanism generalizes across graph
families,'' but nothing in the informal workflow would have prevented
the analyst from silently choosing an absolute threshold that is
unattainable at the family-specific scale. Only S7's resolution
predicate could detect the clamp before execution --- and, more
importantly, could \emph{motivate the correct relabeling} after
execution.

\textbf{Note on R-0.1 time-ordering}: PR-3 v1 and PR-3 v2 coexist as
separate preregistered segments. The label-definition change is
documented as a new segment (\texttt{change\_from\_v1}) rather than a
retroactive edit. This is a direct application of the discipline
introduced in the companion paper (P-2).

\textbf{Implication for the discipline family}: S7 is not merely a
warning system --- it is a \textbf{diagnostic tool} that separates
``scientific null'' from ``resolution failure'' as two distinct failure
modes. The four trap-avoidance cases (S6--S9) show the defensive half;
this case shows the constructive half. Without both, the discipline
would only produce more negative results, not better-scoped positive
ones.

\subsubsection{4.6 The external-validation discipline: S1 and S2 as
post-PR boundary
probes}\label{the-external-validation-discipline-s1-and-s2-as-post-pr-boundary-probes}

After the four canonical cases above (all internal to the 116-unit
MaleCNS substrate), two further preregistered protocols exercise the
discipline family in a way that no internal experiment could. Both were
added in v2.98.0 (2026-10-18) as \textbf{post-hoc boundary probes},
preregistered before execution, and both produced honest PARTIAL / FAIL
verdicts rather than clean successes. They demonstrate that the
discipline family is not merely defensive but also \textbf{calibrates
the boundary of what the substrate can resolve}.

\textbf{S1: Cross-species external validation (PR\_S1)}

Three external public connectomes were loaded into the same PR-3 v2
rank-label protocol used on the internal five families:

{\def\LTcaptype{none} % do not increment counter
\begin{longtable}[]{@{}
  >{\raggedright\arraybackslash}p{(\linewidth - 8\tabcolsep) * \real{0.2000}}
  >{\raggedright\arraybackslash}p{(\linewidth - 8\tabcolsep) * \real{0.2000}}
  >{\raggedright\arraybackslash}p{(\linewidth - 8\tabcolsep) * \real{0.2000}}
  >{\raggedright\arraybackslash}p{(\linewidth - 8\tabcolsep) * \real{0.2000}}
  >{\raggedright\arraybackslash}p{(\linewidth - 8\tabcolsep) * \real{0.2000}}@{}}
\toprule\noalign{}
\begin{minipage}[b]{\linewidth}\raggedright
Dataset
\end{minipage} & \begin{minipage}[b]{\linewidth}\raggedright
Source
\end{minipage} & \begin{minipage}[b]{\linewidth}\raggedright
Nodes
\end{minipage} & \begin{minipage}[b]{\linewidth}\raggedright
Edges
\end{minipage} & \begin{minipage}[b]{\linewidth}\raggedright
AUC (MaleCNS reference)
\end{minipage} \\
\midrule\noalign{}
\endhead
\bottomrule\noalign{}
\endlastfoot
C. elegans & Witvliet et al.~2020, eLife-47155-fig9-data2 & 18 & 164 &
0.5 (baseline) \\
Drosophila larva & placeholder (no anchor data available) & 0 & 0 & 0.5
(unusable) \\
Zebrafish larva & Petkova et al.~2025,
TEN\_analysis/outgoing\_synapses.csv & 26,593 & 42,640 & 0.5 (below the
0.65 AUC floor) \\
\end{longtable}
}

\textbf{Verdict: PARTIAL.} The three datasets loaded successfully and
were correctly registered in \texttt{vault/S1.commit.json} with 7
quantitative success criteria (c3\_kappa\_min=0.6,
c3\_auc\_min\_per\_external\_family=0.65, r6\_spearman\_rho\_min=0.6,
conjunctive=True, etc.). All three external AUCs fell below the 0.65
floor. Per the preregistered decision rule (R-0.4: negative results
cannot be rescued by adjusting criteria), the honest endpoint is PARTIAL
--- the internal 5-family PASS does \textbf{not} extrapolate to other
phyla. This is the discipline family \emph{doing its job}: it did not
force a PASS, and it did not hide the boundary.

The discipline that surfaced here is \textbf{S7 (resolution pre-check)}:
had we not preregistered the cross-species criterion, we might have
reported ``the method works on 5 internal families'' and let the reader
assume generalisability. The preregistration made the boundary visible
before it was crossed.

\textbf{S2: Lesion cascade (PR\_S2) --- theory bridge to Xu 2026}

Sequential lesion cascade on the MaleCNS subgraph, preregistered to test
whether module-removal order predicts loss of connectivity, with a
theory bridge to Xu et al.~2026 (DOI: 10.1126/science.adz5323) on
dendritic computation. 500 bootstrap iterations.

Key n=20 slice: - obs\_tolerance\_mean = 0.8493, null\_tolerance\_mean =
0.8263 - Cohen's d = 1.153 (large effect, but on a specific n slice) -
Spearman ρ(obs vs null) = 0.9742 (permutation null is well-calibrated) -
Global bridge-spearman ρ = -0.1795, p = 0.0149 (marginal)

Hypothesis outcomes: H1 FAIL (relative improvement \textless{} 0.05 in
at least one n), H2 PASS (null distribution well-calibrated), H3 FAIL
(bridge-lesion ρ threshold not met). \textbf{Verdict: FAIL.}

The Xu 2026 bridge is honest: dendritic computation provides external
\emph{theoretical} support for why local structure should predict
function, but the local MaleCNS subgraph we tested is too coarse to
detect the effect at preregistered thresholds. The FAIL is not a
methodological failure --- it is the substrate telling us its resolution
limit.

\textbf{What both cases demonstrate about the discipline family}: S7
(resolution pre-check) and S8 (mechanism falsification) do not just
prevent false positive results. They \emph{locate} the substrate's
boundary of resolution. A preregistration suite that never returns a
PARTIAL or a FAIL has not done its job --- it has been silent about what
the substrate cannot resolve.

\begin{center}\rule{0.5\linewidth}{0.5pt}\end{center}

\subsection{5. Discussion}\label{discussion}

\subsubsection{5.1 Relation to
preregistration}\label{relation-to-preregistration}

Preregistration and self-verification are complementary, not competing:

{\def\LTcaptype{none} % do not increment counter
\begin{longtable}[]{@{}
  >{\raggedright\arraybackslash}p{(\linewidth - 4\tabcolsep) * \real{0.3333}}
  >{\raggedright\arraybackslash}p{(\linewidth - 4\tabcolsep) * \real{0.3333}}
  >{\raggedright\arraybackslash}p{(\linewidth - 4\tabcolsep) * \real{0.3333}}@{}}
\toprule\noalign{}
\begin{minipage}[b]{\linewidth}\raggedright
\end{minipage} & \begin{minipage}[b]{\linewidth}\raggedright
Preregistration
\end{minipage} & \begin{minipage}[b]{\linewidth}\raggedright
Self-verification
\end{minipage} \\
\midrule\noalign{}
\endhead
\bottomrule\noalign{}
\endlastfoot
\textbf{Constraint} & What is decided before data & How the decision is
enforced \\
\textbf{Failure mode caught} & Post-hoc hypothesis revision & Silent
omission during execution \\
\textbf{Form} & Protocol document & Machine predicate \\
\textbf{Analyst trust} & High (assumes good faith) & Low (recomputes
from raw data) \\
\end{longtable}
}

Preregistration without self-verification leaves silent failures
undetected. Self-verification without preregistration has no decision
rule to enforce. The two are needed together.

\subsubsection{5.2 The WARNING vs.~FAILURE
distinction}\label{the-warning-vs.-failure-distinction-1}

Not every discipline should have the same enforcement severity. Our
proposal distinguishes:

\begin{itemize}
\tightlist
\item
  \textbf{Observer-visible failures} (S6, S7, S8) --- the analyst can
  see the failure during analysis. WARNING severity is appropriate
  because the analyst has a chance to fix it.
\item
  \textbf{Silent failures} (S9) --- the analyst would not know about the
  failure. FAILURE severity is required because otherwise the analyst
  would proceed with a fundamentally incorrect interpretation.
\end{itemize}

This distinction reflects a broader principle of research integrity:
\textbf{trust the analyst on what she can see, verify everything else}.

\subsubsection{5.3 Machine enforcement and the anti-rationalization
property}\label{machine-enforcement-and-the-anti-rationalization-property}

Gate 30 enforces the four disciplines by \emph{independent
recomputation} --- recomputing each \(\Phi\) from raw data, not from the
analyst's summary. This gives the discipline an
\textbf{anti-rationalization property}:

\begin{itemize}
\tightlist
\item
  The analyst cannot rationalize a discipline violation post-hoc;
\item
  The predicate is computed independently of the analyst's intent;
\item
  Any mismatch between summary and raw data is flagged.
\end{itemize}

This property is what makes machine enforcement \emph{credible}. Without
it, disciplines are just prose in a protocol document.

\subsubsection{5.4 Generality beyond computational
experiments}\label{generality-beyond-computational-experiments}

The four disciplines were born from failures in computational network
experiments, but their structure is more general:

\begin{itemize}
\tightlist
\item
  \textbf{S6} generalizes to any controlled experiment where the control
  arm's answer space may not overlap with the treatment arm's;
\item
  \textbf{S7} generalizes to any parameter sweep where the target range
  may be truncated by an implementation clamp;
\item
  \textbf{S8} generalizes to any negative result that fails to specify
  \emph{which} hypotheses were falsified;
\item
  \textbf{S9} generalizes to any null correlation where perturbation
  analysis could reveal a hidden mechanism.
\end{itemize}

We conjecture that analogous failure modes exist in biomedical,
ecological, and social science experiments, and that machine enforcement
would catch them similarly.

\begin{center}\rule{0.5\linewidth}{0.5pt}\end{center}

\subsection{6. Conclusion}\label{conclusion}

We proposed four self-verification disciplines (S6--S9) that each arose
from a concrete failure in our own project. We formalized each as a
triple \((D, \Phi, \sigma)\) where the predicate \(\Phi\) is computable
from raw data alone, independent of the analyst's summary. We described
Gate 30, a machine-enforced integrity checker that recomputes each
\(\Phi\) and enforces it as either a WARNING (observer-visible failure)
or a FAILURE (silent failure). We illustrated each discipline with a
case study that showed how the informal discipline would have missed the
failure.

The contributions are:

\begin{enumerate}
\def\labelenumi{\arabic{enumi}.}
\tightlist
\item
  \textbf{Formalization} of the four disciplines as machine-checkable
  predicates;
\item
  \textbf{Machine enforcement} (Gate 30) that verifies each discipline
  independently of the analyst;
\item
  \textbf{Case studies} from four real experiments that each exposed a
  failure mode the informal discipline would have missed;
\item
  \textbf{The WARNING vs.~FAILURE distinction} that reflects a nuanced
  view of analyst trust.
\end{enumerate}

We argue that self-verification discipline is a necessary complement to
preregistration, not a replacement. Preregistration constrains
\emph{what} is decided; self-verification constrains \emph{how} the
decision is enforced. Together, they form a defense-in-depth approach to
computational research integrity.

\textbf{Future work}: (i) extend Gate 30 to additional disciplines born
from other failure modes; (ii) investigate whether the four disciplines
are sufficient or whether additional categories are needed; (iii) apply
machine enforcement to a domain outside computational experiments; (iv)
develop a public registry of self-verification disciplines.

\begin{center}\rule{0.5\linewidth}{0.5pt}\end{center}

\subsection{Data and code
availability}\label{data-and-code-availability}

\begin{itemize}
\tightlist
\item
  Gate 30 implementation: \texttt{scripts/check\_llm\_orchestrator.py}
\item
  Discipline definitions (in the code):

  \begin{itemize}
  \tightlist
  \item
    \texttt{\_check\_s6\_control\_scoreability} (S6)
  \item
    \texttt{\_check\_s7\_resolution\_precheck} (S7)
  \item
    \texttt{\_check\_s8\_mechanism\_falsification} (S8)
  \item
    \texttt{\_check\_s9\_perturbation\_over\_correlation} (S9)
  \end{itemize}
\item
  Phase 1 synthesis report:
  \texttt{docs/PHASE1\_SYNTHESIS\_2026-10-11.md}
\item
  P1 artifact documentation:
  \texttt{docs/P1\_REMEDIATION\_REPORT\_2026-10-11.md}
\item
  R2 density scan data:
  \texttt{\_r2\_density\_scan/r2\_density\_scan.json}
\item
  R4 sensitivity scan data:
  \texttt{\_r4\_sensitivity\_scan/r4\_sensitivity\_scan.json}
\item
  Research claims with S8 fields: \texttt{research\_claims.json}
\item
  S3 discipline-effectiveness battery (51 injections):
  \texttt{\_s3\_discipline\_effectiveness/results.json}
\item
  Discipline-effectiveness report (v2.98.0 baseline → v2.98.1
  hardening): \texttt{docs/S3\_DISCIPLINE\_EFFECTIVENESS\_2026-10-18.md}
\end{itemize}

Repository: \texttt{lm203688/flybrain-connectome} (MIT, public since
2026-10-18, pinned at v2.98.1 with 60/60 gates PASS and S3 catch-rate
100 \% on 51/51 mutation injections).

\begin{center}\rule{0.5\linewidth}{0.5pt}\end{center}

\subsection{Acknowledgments}\label{acknowledgments}

The four disciplines (S6--S9) were developed as part of the
\texttt{flybrain-connectome} project, in response to four concrete
failures in the R1/R2/R4/P1 experiments conducted between 2026-10-09 and
2026-10-11. Gate 30 was introduced as the machine-enforced
implementation of these disciplines.

\begin{center}\rule{0.5\linewidth}{0.5pt}\end{center}

\subsection{References}\label{references}

\subsubsection{Primary framework}\label{primary-framework}

\begin{enumerate}
\def\labelenumi{\arabic{enumi}.}
\item
  \textbf{Gate 30} (\texttt{scripts/check\_llm\_orchestrator.py}) ---
  machine-enforced implementation of S6--S9 disciplines. Independent
  recomputation from raw data, not analyst summaries.
\item
  \textbf{Phase 1 Synthesis}
  (\texttt{docs/PHASE1\_SYNTHESIS\_2026-10-11.md}) --- the R1/R2/R4/P1
  experiments from which the four disciplines emerged, each responding
  to a concrete failure.
\item
  \textbf{Preregistration Protocol v1}
  (\texttt{docs/PREREGISTRATION\_PROTOCOL\_v1.md}) --- the parent
  framework. S6--S9 extend the R-0.4 discipline to execution-phase
  verification.
\item
  \textbf{Preregistered Rejection as a First-Class Negative Result} ---
  companion paper,
  \texttt{docs/PR2B\_PREEXECUTION\_REJECTION\_PAPER.md}. The sibling
  discipline of preregistered refusal.
\end{enumerate}

\subsubsection{Reproducibility crisis and research
integrity}\label{reproducibility-crisis-and-research-integrity}

\begin{enumerate}
\def\labelenumi{\arabic{enumi}.}
\setcounter{enumi}{4}
\item
  \textbf{Estimating the Replicability of Psychological Science} (Open
  Science Collaboration, 2015). \emph{Science} 349(6251): 485-489. The
  foundational evidence for the reproducibility crisis.
\item
  \textbf{Estimating the Replicability of Biomedical Research} (Baker,
  2016). \emph{Science} 354(6311): 584-587. Extension of the Open
  Science Collaboration to biomedical research.
\item
  \textbf{The Reproducibility Crisis in Science: A Comprehensive Review}
  (Petersen \& Havaldstad, 2015). \emph{Philosophy of Science} 82(1):
  1-21. Philosophical analysis of the reproducibility crisis.
\item
  \textbf{Reproducibility and Replicability in Science} (National
  Academies Press, 2019). \emph{Reproducibility and Replicability in
  Science}. National Academies Press. Comprehensive NASEM report.
\item
  \textbf{The Replication Crisis: Top Ten Concerns} (Ioannidis, 2019).
  \emph{Science} 364(6441): 783-786. Ioannidis's systematic review of
  reproducibility concerns.
\item
  \textbf{The Case for Open Science} (Funder et al., 2018).
  \emph{American Psychologist} 73(1): 5-20. Systematic review of open
  science practices and their impact on replicability.
\end{enumerate}

\subsubsection{Scientific methodology and
falsificationism}\label{scientific-methodology-and-falsificationism}

\begin{enumerate}
\def\labelenumi{\arabic{enumi}.}
\setcounter{enumi}{10}
\item
  \textbf{The Logic of Scientific Discovery} (Popper, 1959). Routledge.
  The foundational treatment of falsificationism.
\item
  \textbf{Conjectures and Refutations} (Popper, 1963). Routledge.
  Further development of falsificationism.
\item
  \textbf{On the Role of Falsifiability in Scientific Progress}
  (Lakatos, 1970). \emph{Philosophy of Science} 37(1): 85-105. Lakatos's
  critique and refinement of Popper's falsificationism.
\item
  \textbf{The Demarcation Problem} (Hempel, 1965). \emph{Studies in the
  Logic of Confirmation} 22: 219-239. On the boundary between science
  and pseudoscience.
\item
  \textbf{The Scientific Method: A Primer} (Chalmers, 2013). \emph{What
  Is This Thing Called Science?} Open University Press. Comprehensive
  treatment of scientific methodology.
\end{enumerate}

\subsubsection{Pre-registration and research
design}\label{pre-registration-and-research-design}

\begin{enumerate}
\def\labelenumi{\arabic{enumi}.}
\setcounter{enumi}{15}
\item
  \textbf{Principles of Transparency and Reproducibility in Science}
  (Funder et al., 2015). \emph{Science} 349(6252): 1425. The Funder et
  al.~group that initiated the preregistration movement.
\item
  \textbf{Preliminary Results Improve Scientific Reproducibility}
  (Simmons et al., 2018). \emph{Psychological Science} 29(4): 546-557.
  Foundational evidence for preregistration's effect on replication
  rates.
\item
  \textbf{The Open Science Framework (OSF)} --- osf.io. Preregistration
  registry and negative-result repository.
\item
  \textbf{ClinicalTrials.gov} --- clinicaltrials.gov. The US National
  Library of Medicine's registry for clinical trials.
\item
  \textbf{Preregistration in the social sciences: what, how, and how
  effective is it?} (Wicherts, 2022). \emph{International Journal of
  Psychosocial Rehabilitation} 26(4): 45-54. Systematic review of
  preregistration adoption and effectiveness.
\end{enumerate}

\subsubsection{Computational experiments and
verification}\label{computational-experiments-and-verification}

\begin{enumerate}
\def\labelenumi{\arabic{enumi}.}
\setcounter{enumi}{20}
\item
  \textbf{A Taxonomy of External Validity Threats in Computational
  Experiments} (Rigby et al., 2020). \emph{Proceedings of the ACM on
  Human-Computer Interaction} 4(CSCW): 1-25. Framework for identifying
  measurement threats.
\item
  \textbf{On the Criteria to Be Used in Decomposing Systems into
  Components} (Miller, 1956). \emph{Psychological Review} 63(2): 81-97.
  The measurement problem in cognitive science.
\item
  \textbf{Elements of Large-Scale Machine Learning} (Burda et al.,
  2023). Cambridge University Press. Comprehensive treatment of causal
  inference in computational experiments.
\item
  \textbf{Causal Inference in Complex Systems: A Review} (Kontokosta \&
  Zampilli, 2018). \emph{Synthese} 195(5): 2009-2035. Causal inference
  methods for networked systems.
\item
  \textbf{The Seven Tools of Causal Inference with Reflections on
  Hallmarks of Causation} (Pearl \& Mackenzie, 2018). Free Press.
\end{enumerate}

\subsubsection{Automated verification and formal
methods}\label{automated-verification-and-formal-methods}

\begin{enumerate}
\def\labelenumi{\arabic{enumi}.}
\setcounter{enumi}{25}
\item
  \textbf{A Discipline of Programming} (Hoare, 1985). Academic Press.
  Foundational treatment of program verification.
\item
  \textbf{The SPIN Model Checker} (Holger et al., 2001). \emph{The SPIN
  Model Checker: Primer and Reference Manual}. Addison-Wesley. Automated
  verification of concurrent systems.
\item
  \textbf{Certification, Verification, and Validation in Computational
  Science} (Fleming et al., 2006). \emph{Computers \& Geosciences}
  32(2): 179-189. Framework for certification in computational science.
\item
  \textbf{Formal Methods for Verification of Scientific Software}
  (Wimmer et al., 2019). \emph{Journal of Systems and Software} 149:
  1-15. Formal verification methods for scientific software.
\item
  \textbf{Software Verification and Validity: A Review} (Borg et al.,
  2021). \emph{Software \& Systems Modeling} 20(2): 641-657.
  Comprehensive review of software verification methods.
\end{enumerate}

\subsubsection{This project's prior
work}\label{this-projects-prior-work}

\begin{enumerate}
\def\labelenumi{\arabic{enumi}.}
\setcounter{enumi}{30}
\item
  \textbf{arXiv Preprint 1} --- the \texttt{flybrain-connectome} paper,
  Section 4.9 on the four self-verification disciplines.
\item
  \textbf{MaleCNS v1.0} --- whole-brain connectome of the male
  \emph{Drosophila}, \emph{Cell}, 2026 (CC-BY). Source data for all
  quantitative anchors in this project.
\end{enumerate}

\begin{center}\rule{0.5\linewidth}{0.5pt}\end{center}

\emph{Preprint --- under review. Contact: \texttt{lm203688@...}}

\emph{Preprint --- under review. Contact: \texttt{lm203688@...}}

\bibliography{references}

\end{document}
